\section{The $G_{\alpha\rho}$-Gate: Empirical Falsifiability and $\beta$-Metric}

\subsection{Theoretical Formulation}
We introduce the $G_{\alpha\rho}$-Gate as a formal mechanism to bridge the EABC-model predictions with observable astrophysical signatures. We assume a scalar-field-dependent coupling of the fine-structure constant $\alpha$ to the ambient vacuum density $\rho$. The effective coupling is defined by:
\begin{equation}
    \alpha(\rho) = \alpha_0 \, [1 + \beta \cdot f(\rho)]
\end{equation}
where $\alpha_0$ represents the laboratory-measured value, $\beta \simeq 7.9 \times 10^{-3}$ is the derived model constant, and $f(\rho)$ denotes the environmental screening function. We impose the following boundary conditions:
\begin{itemize}
    \item \textbf{Local Screening:} $f(\rho_{\text{local}}) = 0$, ensuring consistency with terrestrial high-precision experiments and atomic clock constraints.
    \item \textbf{Asymptotic Limit:} $f(\rho \to 0) = 1$, defining the "Zero-Density-Limit" in the intergalactic medium (IGM).
\end{itemize}

\subsection{Falsification Criteria ($A\rho$-Matrices)}
The EABC-model predicts a relative shift in spectral lines, $\Delta \omega / \omega$, for metallic absorbers in low-density environments. Utilizing the Many-Multiplet (MM) method, the predicted sensitivity is governed by:
\begin{equation}
    \frac{\Delta \omega}{\omega} \simeq K_i \cdot \beta \cdot f(\rho_{\text{void}})
\end{equation}
where $K_i$ are the ionic sensitivity coefficients. Adopting $f(\rho_{\text{void}}) \approx 1$, the model yields the following predictions for selected ionic species:

\begin{table}[h!]
\centering
\begin{tabular}{|l|c|c|}
\hline
\textbf{Ion} & \textbf{$K_i$} & \textbf{$\Delta \omega / \omega$ (Pred.)} \\ \hline
Mg II        & 0.05           & $+3.95 \times 10^{-4}$                 \\ \hline
Fe II        & 0.09           & $+7.11 \times 10^{-4}$                 \\ \hline
Cr II        & 0.15           & $+1.19 \times 10^{-3}$                 \\ \hline
Zn II        & 0.22           & $+1.74 \times 10^{-3}$                 \\ \hline
\end{tabular}
\caption{Predicted spectral shifts for selected ions under the EABC-model assumption.}
\end{table}

\subsection{Discussion of the "Claim-Wall"}
Current high-resolution spectroscopic observations (e.g., VLT/ESPRESSO) constrain relative variations of $\alpha$ to the order of $\sim 10^{-6}$. The EABC-model prediction of $\sim 10^{-3}$ significantly exceeds these observational limits. 

We explicitly acknowledge this discrepancy not as a model failure, but as a robust Falsification-Gate. The EABC-model is thus uniquely defined by a "Claim-Wall": If subsequent deep-space spectroscopic analysis of low-density absorbers does not reveal a correlated shift in the specified range, the hypothesis of a universal $\alpha$-coupling to the local vacuum density must be rejected. This operational approach shifts the discourse from vague phenomenological speculation to a hard, testable quantitative framework, inviting empirical verification or immediate refutation within the standard cosmological observational paradigm.